В этом разделе мы приведем описание и результаты работы алгоритма на нескольких классических моделях
\subsection{Gaussian Mixture}
\subsection{SIR}
\hfill \break
\hfill \break
Здесь мы рассмотрим модель распространения заболевания SIR (Suspectable-Infected-Recovered)
Она устроена следующим образом. В каждый момент времени состояние описывается тремя числами $S, I, R$ 
- которые означают число людей из соответстующей категории ($S$ - Здоровые, $I$ - Больные, $R$ - Переболевшие). 
Всегда выполнено соотношение $S + I + R = N$ - общее число людей, а динамика задается так
\[
\begin{cases}
S_{t+1} = S_t - \Delta_{inf} \\
I_{t+1} = I_t + \Delta_{inf} - \Delta_{rec} \\
R_{t+1} = R_t + \Delta_{rec} \\
\end{cases}
\]
где $\Delta_{inf} \sim Binomial(S, 1 - e^{-\beta \frac{I}{N}}), \Delta_{rec} \sim Binomial(I, 1 - e^{-\gamma})$. 
Здесь $\beta, \gamma$ --- искомые параметры $\theta \in \mathbb{R}^2$, а $x \in \mathbb{R}^5$ --- вектор саммари-статистик, 
полученный из временного ряда $(S_t, I_t, R_t) \in \mathbb{R}^{T \times 3}$ 


\subsection{Модель Хестона}
\hfill \break
\hfill \break
Модель Хестона стохастической волатильности описывается стохастическим дифференциальным уравнением
\[
blah-blah
\]
В рамках задачи, мы воспринимаем его как следующую процедуру генерации
% \begin{algorithm}
% \end{algorithm}

алгоритм

Искомыми параметрами здесь являются $\mu$ --- , $\sigma$ ---, $\phi$ ---, $m$ --- . 
В качестве наблюдений мы используем выход последнего скрытого слоя реккурентной нейросети, 
запущенной на временном ряде цены актива. Эта сеть обучается вместе с векторным полем. 
